\documentclass[a4paper]{article}

% \usepackage[default]{fontsetup}

\usepackage{fancyhdr}
\usepackage{extramarks}
\usepackage{amsmath}
\usepackage{amsthm}
\usepackage{amsfonts}
\usepackage{tikz}
\usepackage[plain]{algorithm}
\usepackage{algpseudocode}
\usepackage{enumerate}
\usepackage{tikz}

\usetikzlibrary{automata,positioning}

%
% Basic Document Settings
%  

\topmargin=-0.2in
\evensidemargin=0in
\oddsidemargin=0in
\textwidth=6.5in
\textheight=9.5in
\headsep=0.25in

\linespread{1.1}

\pagestyle{fancy}
\lhead{\hmwkAuthorName}
\chead{\hmwkClass : \hmwkTitle}
\rhead{\firstxmark}
\lfoot{\lastxmark}
\cfoot{\thepage}

\renewcommand\headrulewidth{0.4pt}
\renewcommand\footrulewidth{0.4pt}

\setlength\parindent{0pt}

%
% Create Problem Sections
%

\newcommand{\enterProblemHeader}[1]{
    \nobreak\extramarks{}{Problem \arabic{#1} continued on next page\ldots}\nobreak{}
    \nobreak\extramarks{Problem \arabic{#1} (continued)}{Problem \arabic{#1} continued on next page\ldots}\nobreak{}
}

\newcommand{\exitProblemHeader}[1]{
    \nobreak\extramarks{Problem \arabic{#1} (continued)}{Problem \arabic{#1} continued on next page\ldots}\nobreak{}
    \stepcounter{#1}
    \nobreak\extramarks{Problem \arabic{#1}}{}\nobreak{}
}

\newcommand*\circled[1]{\tikz[baseline=(char.base)]{
		\node[shape=circle,draw,inner sep=2pt] (char) {#1};}}


\setcounter{secnumdepth}{0}
\newcounter{partCounter}
\newcounter{homeworkProblemCounter}
\setcounter{homeworkProblemCounter}{1}
\nobreak\extramarks{Problem \arabic{homeworkProblemCounter}}{}\nobreak{}

%
% Homework Problem Environment
%
% This environment takes an optional argument. When given, it will adjust the
% problem counter. This is useful for when the problems given for your
% assignment aren't sequential. See the last 3 problems of this template for an
% example.
%

\newenvironment{homeworkProblem}[2][-1]{%
    \ifnum#1>0\relax
        \setcounter{homeworkProblemCounter}{#1}
    \fi
    \section{Problem \arabic{homeworkProblemCounter} (#2 points)}
    \setcounter{partCounter}{1}
    \enterProblemHeader{homeworkProblemCounter}
}{%
    \exitProblemHeader{homeworkProblemCounter}
}

%
% Homework Details
%   - Title
%   - Class
%   - Due date
%   - Name
%   - Student ID

\newcommand{\hmwkTitle}{Homework\ \#06}
\newcommand{\hmwkClass}{Probability \& Statistics for EECS}
\newcommand{\hmwkDueDate}{2026-11-1}
\newcommand{\hmwkAuthorName}{Zhang Yiheng}
\newcommand{\hmwkAuthorID}{2025531068}


%
% Title Page
%

\title{
    \vspace{2in}
    \textmd{\textbf{\hmwkClass:\\  \hmwkTitle}}\\
    \normalsize\vspace{0.1in}\small{Due\ on\ \hmwkDueDate\ at 23:59}\\
	\vspace{4in}
}

\author{
	Name: \textbf{\hmwkAuthorName} \\
	Student ID: \hmwkAuthorID}
\date{}

\renewcommand{\part}[1]{\textbf{\large Part \Alph{partCounter}}\stepcounter{partCounter}\\}

%
% Various Helper Commands
%

% Useful for algorithms
\newcommand{\alg}[1]{\textsc{\bfseries \footnotesize #1}}
% For derivatives
\newcommand{\deriv}[1]{\frac{\mathrm{d}}{\mathrm{d}x} (#1)}
% For partial derivatives
\newcommand{\pderiv}[2]{\frac{\partial}{\partial #1} (#2)}
% Integral dx
\newcommand{\dx}{\mathrm{d}x}
% Alias for the Solution section header
\newcommand{\solution}{\textbf{\large Solution}}
% Probability commands: Expectation, Variance, Covariance, Bias
\newcommand{\E}{\mathrm{E}}
\newcommand{\Var}{\mathrm{Var}}
\newcommand{\Cov}{\mathrm{Cov}}
\newcommand{\Bias}{\mathrm{Bias}}

\begin{document}


% \maketitle
% \thispagestyle{empty}
% \pagebreak

\date{
Due on Nov. 1, 2026, 23:59 UTC+8}
\title{SI 140A.02  Probability \& Statistics for EECS, Fall 2026 \\
Homework 6}
\maketitle
Read all the instructions below carefully before you start working on the assignment, and before you make a submission.
\begin{itemize}
    \item You are required to submit your homework in PDF format generated by \textbf{LaTeX}, otherwise you will get no points of this homework.
    \item You are required to write down all the major steps towards making your conclusions; otherwise you may obtain limited points of the problem.
    \item Write your homework in English; otherwise you will get no points of this homework.
    \item Any form of plagiarism will lead to $0$ point of this homework. 
\end{itemize}
\newpage

\begin{homeworkProblem}{20}

Let $n\ge 3$, and let $X_1,\ldots,X_n$ be binary random variables
with $\Pr(X_i=1)=p_i\in(0,1)$. Suppose that every collection of
$n-1$ variables is mutually independent. No symmetry of the
joint distribution is assumed.

\begin{enumerate}
\item[(a)] Determine all possible joint PMFs and the exact
admissible range of their free parameter. Prove both necessity
and sufficiency. 
\item[(b)] Suppose additionally that $n=4$, $p_i=1/2$ for all $i$,
and all $16$ joint outcomes have strictly positive probability.
Determine the exact range of
$\Pr(X_1+X_2+X_3+X_4=2)$.
Does the value $3/8$ force mutual independence?
\end{enumerate}

\solution

\medskip
\textbf{(a)} Write $q_i=1-p_i$. For $x=(x_1,\ldots,x_n)\in\{0,1\}^n$,
let
\[
 b(x)=\prod_{i=1}^n p_i^{x_i}q_i^{1-x_i},
 \qquad z(x)=\#\{i:x_i=0\}.
\]
We claim that all possible joint PMFs have the form
\[
 \boxed{\Pr(X_1=x_1,\ldots,X_n=x_n)=b(x)+(-1)^{z(x)}c,}
\]
where the exact admissible range is
\[
 \boxed{-\min_{x:\,z(x)\text{ even}}b(x)
 \ \le c\le\ \min_{x:\,z(x)\text{ odd}}b(x).}
\]

\textit{Necessity.}
Set $c=\Pr(X_1=\cdots=X_n=1)-\prod_{i=1}^n p_i$.
For any index $j$, summing over the two values of $X_j$ and using
the independence of the other $n-1$ variables gives
\[
 \Pr(X_i=x_i\text{ for all }i\ne j)
 =\prod_{i\ne j}p_i^{x_i}q_i^{1-x_i}.
\]
In particular, an outcome with its only zero at coordinate $j$ has probability
\[
 \prod_{i\ne j}p_i
 -\left(\prod_{i=1}^n p_i+c\right)
 =q_j\prod_{i\ne j}p_i-c.
\]
Now induct on the number of zeros. The case of no zeros follows from
the definition of $c$. Suppose the formula holds for every outcome
with $t-1$ zeros. For an outcome $x$ with $t$ zeros, choose $j$ with
$x_j=0$, and let $x^{(j)}$ be obtained by changing that zero to one.
Writing $B=\prod_{i\ne j}p_i^{x_i}q_i^{1-x_i}$, we obtain
\begin{align*}
 \Pr(X=x)
 &=B-\Pr(X=x^{(j)})\\
 &=B-\bigl(p_jB+(-1)^{t-1}c\bigr)\\
 &=q_jB+(-1)^tc=b(x)+(-1)^{z(x)}c.
\end{align*}
Nonnegativity requires $c\ge-b(x)$ for every even $z(x)$ and
$c\le b(x)$ for every odd $z(x)$, giving exactly the stated bounds.

\textit{Sufficiency.}
Conversely, for any $c$ in this interval, all the proposed probabilities
are nonnegative. Their sum is one because
\[
 \sum_{x\in\{0,1\}^n}b(x)=\prod_{i=1}^n(p_i+q_i)=1,
 \qquad
 \sum_{x\in\{0,1\}^n}(-1)^{z(x)}=(1-1)^n=0.
\]
Moreover, for any fixed values of the coordinates other than $j$,
summing over $x_j\in\{0,1\}$ cancels the two correction terms and gives
\[
 \sum_{x_j=0}^1\Pr(X=x)
 =\prod_{i\ne j}p_i^{x_i}q_i^{1-x_i}.
\]
Thus every collection of $n-1$ variables is mutually independent with
the required marginals. This also proves that both endpoints of the
parameter interval are admissible.

\medskip
\textbf{(b)} Here $b(x)=1/16$ for every outcome. Strict positivity is
equivalent to
\[
 -\frac1{16}<c<\frac1{16}.
\]
Each outcome with exactly two ones also has two zeros, so
\[
 \Pr(X_1+X_2+X_3+X_4=2)
 =\binom42\left(\frac1{16}+c\right)
 =\frac38+6c.
\]
Consequently, the exact range is $\boxed{(0,3/4)}$, with neither
endpoint attained. The value $3/8$ forces $c=0$, in which case the
joint PMF equals the product of the marginals. Hence
\textbf{yes}, this value forces mutual independence under the given assumptions.
\end{homeworkProblem}
\newpage
\begin{homeworkProblem}{20}
A population contains $m$ red balls and $m$ blue balls.
A uniformly chosen sample of size $n$ is drawn without replacement.
Let $X_{m,n}$ be the number of red balls, and let
$Y_n\sim\operatorname{Bin}(n,1/2)$.

\begin{enumerate}
\item[(a)] For fixed positive integer $n$ and fixed
$k\in\{0,\ldots,n\}$, give the PMF of $X_{m,n}$ and evaluate
\[
\lim_{m\to\infty}
m\left(\frac{\Pr(X_{m,n}=k)}{\Pr(Y_n=k)}-1\right).
\]
Do not merely quote convergence of the hypergeometric distribution.

\item[(b)] Set $n=2m-1$ and $A_m=\{m-1,m\}$. Evaluate
\[
\lim_{m\to\infty}
\left[\Pr(X_{m,2m-1}\in A_m)
-\Pr(Y_{2m-1}\in A_m)\right].
\]
Does a growing population with a fixed red proportion alone
justify the binomial approximation?
Prove the required limit without the central limit theorem
or Stirling's formula. 
\end{enumerate}

\solution

\medskip
\textbf{(a)} For $1\le n\le 2m$, counting samples gives
\[
 \Pr(X_{m,n}=k)
 =\frac{\binom{m}{k}\binom{m}{n-k}}{\binom{2m}{n}},
 \qquad \max(0,n-m)\le k\le\min(n,m),
\]
and the probability is zero for other integers $k$.
Fix $n$ and $k\in\{0,\ldots,n\}$. For all sufficiently large $m$,
the ratio to the binomial PMF is
\begin{align*}
 R_m
 &:=\frac{\Pr(X_{m,n}=k)}{\Pr(Y_n=k)}\\
 &=\frac{2^n(m)_k(m)_{n-k}}{(2m)_n}\\
 &=\frac{\displaystyle\prod_{j=0}^{k-1}\left(1-\frac jm\right)
                 \prod_{j=0}^{n-k-1}\left(1-\frac jm\right)}
         {\displaystyle\prod_{j=0}^{n-1}\left(1-\frac j{2m}\right)},
\end{align*}
where $(a)_s=a(a-1)\cdots(a-s+1)$ and empty products equal one.
Taking natural logarithms turns these products into sums:
\[
 \ln R_m
 =\sum_{j=0}^{k-1}\ln\left(1-\frac jm\right)
  +\sum_{j=0}^{n-k-1}\ln\left(1-\frac jm\right)
  -\sum_{j=0}^{n-1}\ln\left(1-\frac j{2m}\right).
\]
Since $n$ and $k$ are fixed, Taylor's expansion
$\ln(1-u)=-u+O(u^2)$ yields
\begin{align*}
 \ln R_m
 &=-\frac{k(k-1)}{2m}
   -\frac{(n-k)(n-k-1)}{2m}
   +\frac{n(n-1)}{4m}+O(m^{-2})\\
 &=\frac{n-(2k-n)^2}{4m}+O(m^{-2}).
\end{align*}
Using $e^v=1+v+O(v^2)$, we therefore have
\[
 R_m=1+\frac{n-(2k-n)^2}{4m}+O(m^{-2}),
 \qquad
 \boxed{\lim_{m\to\infty}m(R_m-1)=\frac{n-(2k-n)^2}{4}.}
\]

\medskip
\textbf{(b)} A sample of size $2m-1$ omits exactly one ball.
It contains $m-1$ red balls if the omitted ball is red and $m$ red
balls otherwise. Thus
\[
 \Pr(X_{m,2m-1}\in A_m)=1.
\]
By symmetry and the identity
$\binom{2m-1}{m}=\tfrac12\binom{2m}{m}$,
\[
 \Pr(Y_{2m-1}\in A_m)
 =\frac{2\binom{2m-1}{m}}{2^{2m-1}}
 =2\frac{\binom{2m}{m}}{4^m}.
\]
To prove that this tends to zero, write the factorial ratio as
\[
 a_m:=\frac{\binom{2m}{m}}{4^m}
 =\prod_{j=1}^m\frac{2j-1}{2j}
 =\prod_{j=1}^m\left(1-\frac1{2j}\right).
\]
The Taylor remainder in $\ln(1-u)=-u+O(u^2)$ is bounded uniformly
for $0\le u\le1/2$. Hence
\[
 \ln a_m
 =-\frac12\sum_{j=1}^m\frac1j
   +O\left(\sum_{j=1}^m\frac1{j^2}\right)
 =-\frac12\sum_{j=1}^m\frac1j+O(1)
 \longrightarrow-\infty.
\]
Here the harmonic series diverges, whereas $\sum_{j=1}^\infty j^{-2}$
converges. Therefore $a_m\to0$, without using the central limit
theorem or Stirling's formula, and
\[
 \boxed{\lim_{m\to\infty}
 \left[\Pr(X_{m,2m-1}\in A_m)-\Pr(Y_{2m-1}\in A_m)\right]=1.}
\]
Thus \textbf{no}: a growing population with a fixed red proportion alone
does not justify the binomial approximation. In this example the sampling
fraction tends to one, and the two probabilities above differ by an
amount tending to one.
\end{homeworkProblem}
\newpage
\begin{homeworkProblem}{20}
Let $X$ take values in $\{1,2,3,4\}$, with zero probabilities
allowed. Suppose
\[
\Pr(X\in\{1,2\})=\frac23,
\qquad
\Pr(X\in\{2,3\})=\frac12.
\]
Entropy is measured in bits, with $0\log_2 0=0$.

\begin{enumerate}
\item[(a)] Determine the maximum and minimum of $H(X)$ and
all PMFs attaining them. Prove that the extrema are global.

\item[(b)] If all four probabilities must be strictly positive,
determine the exact set of possible entropy values and state
which boundaries are attained. 
\end{enumerate}

\solution

\medskip
\textbf{(a)} Set $x=\Pr(X=2)$. The two constraints and normalization give
\[
 (p_1,p_2,p_3,p_4)
 =\left(\frac23-x,\ x,\ \frac12-x,\ x-\frac16\right).
\]
All four probabilities are nonnegative if and only if
$1/6\le x\le1/2$. Thus the entropy is
\begin{align*}
 H(x)={}&-\left(\frac23-x\right)\log_2\left(\frac23-x\right)
         -x\log_2 x\\
       &-\left(\frac12-x\right)\log_2\left(\frac12-x\right)
         -\left(x-\frac16\right)\log_2\left(x-\frac16\right),
 \qquad \frac16\le x\le\frac12,
\end{align*}
with $0\log_2 0=0$. This function is continuous on the closed interval.
For $1/6<x<1/2$, differentiation yields
\[
 H'(x)=\log_2\frac{(\frac23-x)(\frac12-x)}{x(x-\frac16)},
\]
so
\[
 H'(x)=0
 \iff \left(\frac23-x\right)\left(\frac12-x\right)
       =x\left(x-\frac16\right)
 \iff x=\frac13.
\]
Moreover,
\[
 H''(x)=-\frac1{\ln2}
 \left(\frac1{\frac23-x}+\frac1x
       +\frac1{\frac12-x}+\frac1{x-\frac16}\right)<0.
\]
Hence $H$ is strictly concave, and the unique stationary point is the
unique global maximum. It is attained at
\[
 \boxed{(p_1,p_2,p_3,p_4)=\left(\frac13,\frac13,\frac16,\frac16\right),
 \qquad H_{\max}=\log_2 3+\frac13.}
\]
For the minimum, the two endpoint PMFs are
\[
 x=\frac16:\quad\left(\frac12,\frac16,\frac13,0\right),
 \qquad
 x=\frac12:\quad\left(\frac16,\frac12,0,\frac13\right).
\]
They have the same entropy,
\[
 H\left(\frac16\right)=H\left(\frac12\right)
 =-\frac12\log_2\frac12-\frac16\log_2\frac16
  -\frac13\log_2\frac13
 =\frac23+\frac12\log_2 3.
\]
Every interior point is a nontrivial convex combination of the endpoints;
strict concavity makes its entropy strictly greater than their common
value. Therefore
\[
 \boxed{H_{\min}=\frac23+\frac12\log_2 3,}
\]
and the two endpoint PMFs above are all the minimizers.

\medskip
\textbf{(b)} Strict positivity is equivalent to $1/6<x<1/2$.
The maximizing PMF from part (a) is still admissible, but neither
minimizing PMF is. In addition, $H$ is continuous and strictly increasing
on $(1/6,1/3]$, with
\[
 \lim_{x\downarrow1/6}H(x)=\frac23+\frac12\log_2 3,
 \qquad H\left(\frac13\right)=\log_2 3+\frac13.
\]
The intermediate value theorem therefore shows that the exact set of
possible entropy values is
\[
 \boxed{\left(\frac23+\frac12\log_2 3,\ \log_2 3+\frac13\right].}
\]
The upper boundary is attained uniquely at the maximizing PMF in part (a).
The lower boundary is an infimum and is not attained.

\end{homeworkProblem}
\newpage
\begin{homeworkProblem}{20}
Let $N\ge2$ be an integer. A random variable $X$ takes values in
$\{0,1,\ldots,N\}$ and satisfies $\mathbb E[X]=\mu$, where
$0<\mu<N$. No other distributional assumptions are imposed.

\begin{enumerate}
\item[(a)] Determine the minimum and maximum possible variance
and characterize all equality cases. Cover both integer and
noninteger values of $\mu$. 
\item[(b)] Prove that every value between these extrema is
attainable, and provide an explicit PMF construction.

\end{enumerate}

\solution

\medskip
\textbf{(a)} The minimum and maximum variances are
\[
 \boxed{\Var_{\min}=(\mu-\lfloor\mu\rfloor)(\lceil\mu\rceil-\mu),
 \qquad \Var_{\max}=\mu(N-\mu).}
\]

\textit{Minimum.}
There is no integer strictly between $\lfloor\mu\rfloor$ and
$\lceil\mu\rceil$. Thus
$(X-\lfloor\mu\rfloor)(X-\lceil\mu\rceil)\ge0$, and taking expectations gives
\[
 \mathbb E[X^2]\ge
 (\lfloor\mu\rfloor+\lceil\mu\rceil)\mu
 -\lfloor\mu\rfloor\lceil\mu\rceil.
\]
Consequently,
\[
 \Var(X)=\mathbb E[X^2]-\mu^2
 \ge(\mu-\lfloor\mu\rfloor)(\lceil\mu\rceil-\mu).
\]
Equality holds if and only if
$X\in\{\lfloor\mu\rfloor,\lceil\mu\rceil\}$ almost surely.
If $\mu$ is noninteger, the mean constraint uniquely forces
\[
 \boxed{\Pr(X=\lfloor\mu\rfloor)=\lceil\mu\rceil-\mu,
 \qquad \Pr(X=\lceil\mu\rceil)=\mu-\lfloor\mu\rfloor,}
\]
with zero probability elsewhere. If $\mu$ is an integer, the unique
minimizer is $\Pr(X=\mu)=1$, with variance zero.

\textit{Maximum.}
Since $0\le X\le N$, we have $X(N-X)\ge0$, and hence
\[
 \mathbb E[X^2]\le N\mathbb E[X]=N\mu,
 \qquad \Var(X)\le N\mu-\mu^2=\mu(N-\mu).
\]
Equality holds if and only if $X(N-X)=0$ almost surely, or equivalently
$X\in\{0,N\}$ almost surely. The mean constraint uniquely determines
\[
 \boxed{\Pr(X=0)=1-\frac\mu N,\qquad \Pr(X=N)=\frac\mu N,}
\]
with zero probability elsewhere. This characterization holds for both
integer and noninteger $\mu$.

\medskip
\textbf{(b)} \textit{Existence.}
Let $L=(L_j)_{j=0}^N$ and $U=(U_j)_{j=0}^N$ be the minimizing and
maximizing PMFs from part (a), respectively. For $t\in[0,1]$, define
\[
 P_t(j)=(1-t)L_j+tU_j,\qquad j=0,\ldots,N.
\]
These numbers are nonnegative and sum to one, so they define a PMF.
Moreover, every such PMF has the required mean:
\[
 \sum_{j=0}^N jP_t(j)=(1-t)\mu+t\mu=\mu.
\]
Since this mean is fixed, its variance is
\[
 V(t)=\sum_{j=0}^N(j-\mu)^2P_t(j)
 =(1-t)\Var_{\min}+t\Var_{\max}.
\]
Thus $V$ is continuous on $[0,1]$, with $V(0)=\Var_{\min}$ and
$V(1)=\Var_{\max}$. By the intermediate value theorem, every value
between the two extrema is attained by an admissible PMF.

\textit{Explicit construction.}
For a prescribed variance $v\in[\Var_{\min},\Var_{\max}]$, set
\[
 \lambda=\frac{v-(\mu-\lfloor\mu\rfloor)(\lceil\mu\rceil-\mu)}
                  {\mu(N-\mu)-(\mu-\lfloor\mu\rfloor)(\lceil\mu\rceil-\mu)}.
\]
The denominator is strictly positive: the maximizing PMF puts positive
mass at both $0$ and $N$, which cannot both lie in
$\{\lfloor\mu\rfloor,\lceil\mu\rceil\}$ when
$N\ge2$, so it cannot attain the minimum. Hence $0\le\lambda\le1$.
Writing the minimizing PMF explicitly as
\[
 L_j=\begin{cases}
 (\lceil\mu\rceil-\mu)\mathbf{1}_{\{j=\lfloor\mu\rfloor\}}
 +(\mu-\lfloor\mu\rfloor)\mathbf{1}_{\{j=\lceil\mu\rceil\}},
 &\mu\notin\mathbb Z,\\
 \mathbf{1}_{\{j=\mu\}},&\mu\in\mathbb Z,
 \end{cases}
\]
an explicit PMF with variance $v$ is
\[
 \boxed{
 \Pr(X=j)=(1-\lambda)L_j
 +\lambda\left(\left(1-\frac\mu N\right)\mathbf{1}_{\{j=0\}}
                    +\frac\mu N\mathbf{1}_{\{j=N\}}\right),
 \qquad j=0,\ldots,N.
 }
\]
Here $\mathbf{1}_{\{\cdot\}}$ is the indicator; if support points coincide,
their contributions are added. By the preceding argument this PMF has
mean $\mu$, and its variance is
\[
 (1-\lambda)(\mu-\lfloor\mu\rfloor)(\lceil\mu\rceil-\mu)
 +\lambda\mu(N-\mu)=v.
\]

\end{homeworkProblem}
\newpage
\begin{homeworkProblem}{20}
Consider independent Bernoulli trials with success probability
$p\in(0,1)$. Fix an integer $r\ge2$, and let $T_r$ be the
trial number of the $r$th success, including the successful trial.

Fix an integer $m\ge r-1$. An observer is told only that
$T_r>m$, not how many successes occurred during the first $m$
trials. On this event, let $R_m=T_r-m$.

\begin{enumerate}
\item[(a)] Find the complete conditional PMF of $R_m$ given
$T_r>m$. Express it also as a mixture of waiting-time
distributions, specifying the mixture weights.

\item[(b)] Can this conditional distribution be a first-success
distribution with any parameter? For each fixed integer $k\ge1$,
evaluate
\[
\lim_{m\to\infty}\Pr(R_m=k\mid T_r>m)
\]
and identify the limiting distribution.
\end{enumerate}

\solution

\medskip
\textbf{(a)} Write $q=1-p$, and let $S_m$ be the number of successes
in the first $m$ trials. Since $\{T_r>m\}=\{S_m\le r-1\}$, put
\[
 D_m:=\Pr(T_r>m)=\sum_{j=0}^{r-1}\binom mj p^j q^{m-j}.
\]
For $T_r=m+k$, the first $m+k-1$ trials must contain exactly $r-1$
successes, and trial $m+k$ must be a success. Thus the complete
conditional PMF is
\[
 \boxed{f_m(k):=\Pr(R_m=k\mid T_r>m)
 =\frac{\binom{m+k-1}{r-1}p^r q^{m+k-r}}{D_m},
 \qquad k=1,2,\ldots,}
\]
and it is zero otherwise. Here $m\ge r-1$ ensures that every positive
integer $k$ is possible.

For the mixture representation, conditional on $S_m=j$, we still need
$r-j$ successes. Future trials are independent of the first $m$ trials,
so the remaining time has the waiting-time distribution for $r-j$
successes. Its mixture weight is
\[
 w_{m,j}:=\Pr(S_m=j\mid T_r>m)
 =\frac{\binom mj p^j q^{m-j}}{D_m},\qquad j=0,\ldots,r-1.
\]
If $G_s$ denotes the number of trials required for $s$ successes, then
\[
 g_s(k):=\Pr(G_s=k)=
 \begin{cases}
 \binom{k-1}{s-1}p^s q^{k-s},&k=s,s+1,\ldots,\\
 0,&\text{otherwise}.
 \end{cases}
\]
Consequently,
\[
 \boxed{f_m(k)=\sum_{j=0}^{r-1}w_{m,j}\,g_{r-j}(k).}
\]
The weights are nonnegative and sum to one, so this is a mixture of
negative binomial waiting-time distributions.

\medskip
\textbf{(b)} A first-success distribution with parameter $\alpha\in(0,1)$
has PMF $\alpha(1-\alpha)^{k-1}$, whose successive probability ratio is
the constant $1-\alpha$. In contrast, part (a) gives
\[
 \frac{f_m(k+1)}{f_m(k)}
 =q\frac{m+k}{m+k-r+1}
 =q\left(1+\frac{r-1}{m+k-r+1}\right),\qquad k\ge1.
\]
Since $r\ge2$, this ratio depends strictly on $k$. Also $f_m(k)>0$
for every $k\ge1$, excluding the degenerate case $\alpha=1$.
Thus \textbf{no}, for any finite $m$ this is not a first-success
distribution with any parameter.

To find the limit, factor the $j=r-1$ term out of $D_m$. For every
fixed $j<r-1$,
\[
 \frac{\binom mj p^j q^{m-j}}
      {\binom m{r-1}p^{r-1}q^{m-r+1}}
 =\frac{\binom mj}{\binom m{r-1}}
   \left(\frac qp\right)^{r-1-j}\longrightarrow0,
\]
because $\binom mj/\binom m{r-1}=O(m^{j-r+1})$ with $r$ and $j$ fixed.
It follows that
\[
 w_{m,r-1}\longrightarrow1,
 \qquad w_{m,j}\longrightarrow0\quad(j<r-1).
\]
Using the finite mixture in part (a), for every fixed integer $k\ge1$,
\[
 \boxed{\lim_{m\to\infty}\Pr(R_m=k\mid T_r>m)
 =g_1(k)=p(1-p)^{k-1}.}
\]
These limiting probabilities sum to one. Hence the limiting distribution
is the geometric (first-success) distribution with parameter $p$ on
$\{1,2,\ldots\}$.

\end{homeworkProblem}
\end{document}
